\documentclass[a4paper,12pt]{article}
\usepackage[utf8]{inputenc}
\usepackage{amsmath, amssymb}
\usepackage{graphicx}
\usepackage{geometry}
\usepackage{hyperref}
\usepackage{listings}
\usepackage{xcolor}
\usepackage{caption}
\usepackage{tabularx}

\geometry{margin=1in}

% --- Impostazioni per il codice sorgente ---
\lstset{
  basicstyle=\ttfamily\footnotesize,
  keywordstyle=\color{blue}\bfseries,
  commentstyle=\color{gray}\itshape,
  stringstyle=\color{orange},
  numbers=left,
  numberstyle=\tiny,
  breaklines=true,
  frame=single,
  captionpos=b,
  tabsize=4,
  showstringspaces=false
}

\begin{document}

\title{Laboratory Assignment: RISC-V Assembly Programming}
\author{Student: Andrei Stefan}
\date{\today}
\maketitle

\section*{Exercise 1}

Using the same flow described before, write and run a program for calculating the Fibonacci sequence and save it in \texttt{program2.s} (any file name is acceptable, as long as it has the \texttt{.s} extension).

\subsection*{Assembly Program}

\begin{lstlisting}[language={[x86masm]Assembler}, caption={Fibonacci sequence generator in RISC-V Assembly}]
# Data section
.section .data
# Place here your program data. In this example
# two vector of floats, a vector of ints and a single int are defined
T0: .word 0x0BADC0DE

# Code section
.section .text
.globl _start
_start:
# The _start label signals the entry point of your program
# DO NOT CHANGE ITS NAME. It must be "_start", not "start" or "main".

Main:
    # Initialize Fibonacci variables
    li x1, 0      # x1 = a = first Fibonacci number (0)
    li x2, 1      # x2 = b = second Fibonacci number (1)
    li x3, 21     # x3 = count = number of terms to generate
    li x4, 0      # x4 = i = loop counter

    # Loop to generate and print remaining 20 numbers
    addi x4, x4, 1    # i = 1 (start from second iteration)

fib_loop:
    beq x4, x3, End   # if i == count, exit loop

    # Calculate next Fibonacci number
    add x5, x1, x2    # x5 = next = a + b

    # Update variables for next iteration
    mv x1, x2         # a = b
    mv x2, x5         # b = next

    # Increment counter and continue loop
    addi x4, x4, 1
    j fib_loop

End:
    # Exit syscall (needed to end simulation gracefully)
    li a0, 0
    li a7, 93
    ecall
\end{lstlisting}

\subsection*{Equivalent C Code}

\begin{lstlisting}[language=C, caption={Fibonacci sequence generator in C}]
/* Fibonacci sequence generator in C
 * Generates the first 21 Fibonacci numbers using iterative approach
 */
int main() {
    int i;
    int a = 0;
    int b = 1;
    int next;
    int count = 21;

    for (i = 1; i < count; i++) {
        next = a + b;
        a = b;
        b = next;
    }
    return 0;
}
\end{lstlisting}

\newpage
\section*{Exercise 2}

Using the same flow described before, write and run an assembly program called \texttt{program\_3.s} for the RISC-V architecture.

The program must:

\begin{enumerate}
    \item Given two arrays of 10 8-bit integer numbers (\texttt{v1}, \texttt{v2}), check if any element of \texttt{v1} is included in \texttt{v2} at least once. Save the matching values in a third vector (\texttt{v3}).

    \begin{lstlisting}[language={[x86masm]Assembler}]
v1: .byte 2, 6, -3, 11, 9, 18, -13, 16, 5, 1
v2: .byte 4, 2, -13, 3, 9, 9, 7, 16, 4, 7
# Expected output
v3: .byte 2, 9, -13, 16
    \end{lstlisting}

    \item Set three flags (\texttt{flag1}, \texttt{flag2}, \texttt{flag3}) to indicate the following conditions:
    \begin{itemize}
        \item \textbf{Flag 1}: The third vector (\texttt{v3}) is empty.  
        $\Rightarrow$ \texttt{flag1 = 1} if empty, \texttt{0} otherwise.
        \item \textbf{Flag 2}: \texttt{v3} is not empty and elements are strictly increasing (\texttt{v3[i+1] > v3[i]}).  
        $\Rightarrow$ \texttt{flag2 = 1} if true, \texttt{0} otherwise.
        \item \textbf{Flag 3}: \texttt{v3} is not empty and elements are strictly decreasing (\texttt{v3[i+1] < v3[i]}).  
        $\Rightarrow$ \texttt{flag3 = 1} if true, \texttt{0} otherwise.
    \end{itemize}
\end{enumerate}

\subsection*{Explanation of Initial Stalls}

At the beginning of the program, there are a few sequential stalls because the cache memory is empty when execution starts. The \texttt{\_start} section is used to load data into the cache so that it will be available during the execution of the program.

\subsection*{Program Performance}

\begin{table}[h!]
\centering
\caption{Program performance for the specific processor configurations}
\begin{tabularx}{\textwidth}{|X|X|X|X|X|}
\hline
\textbf{Program} & \textbf{Clock cycles} & \textbf{Number of Instructions} & \textbf{CPI} & \textbf{IPC} \\
\hline
program\_1 & 23  & 13  & 1.76923 & 0.56521 \\
program\_2 & 156 & 126 & 1.2381  & 0.8076  \\
program\_3 & 844 & 627 & 1.34609 & 0.7428  \\
\hline
\end{tabularx}
\end{table}

\newpage
\section*{Exercise 3}

Perform execution and CPI measurements of some benchmark programs. The folder contains two additional programs:

\begin{itemize}
    \item \texttt{calculate\_pi.s}
    \item \texttt{insertion\_sort.s}
\end{itemize}

Also, include programs from previous exercises:
\begin{itemize}
    \item \texttt{program\_1.s}
    \item \texttt{program\_2.s}
    \item \texttt{program\_3.s}
\end{itemize}

Assume:
\begin{itemize}
    \item Equal initial weights (20\% each program)
    \item Processor frequency: \textbf{1.75 kHz}
\end{itemize}

Then consider the following scenarios:

\begin{description}
    \item[Scenario 1:] 
    \begin{itemize}
        \item program\_1.s: 1\%
        \item program\_2.s: 50\%
        \item program\_3.s: 13\%
        \item calculate\_pi.s: 25\%
        \item insertion\_sort.s: 11\%
    \end{itemize}

    \item[Scenario 2:] 
    \begin{itemize}
        \item program\_1.s: 10\%
        \item program\_2.s: 5\%
        \item program\_3.s: 50\%
        \item calculate\_pi.s: 10\%
        \item insertion\_sort.s: 25\%
    \end{itemize}

    \item[Scenario 3:] 
    \begin{itemize}
        \item program\_1.s: 20\%
        \item program\_2.s: 30\%
        \item program\_3.s: 1.9\%
        \item calculate\_pi.s: 31.4\%
        \item insertion\_sort.s: 16.7\%
    \end{itemize}
\end{description}

\subsection*{Processor Performance Results}

\begin{table}[h!]
\centering
\caption{Processor performance for different weighted programs}
\begin{tabularx}{\textwidth}{|X|X|X|X|X|}
\hline
\textbf{Program} & \textbf{Initial scenario} & \textbf{Scenario 1} & \textbf{Scenario 2} & \textbf{Scenario 3} \\
\hline
calculate\_pi.s   & 0.46068 & 0.57586 & 0.23034 & 0.72328 \\
insertion\_sort.s  & 1.54388 & 0.85197 & 1.93629 & 1.29344 \\
program\_1.s       & 0.00263 & 0.00013 & 0.00131 & 0.00263 \\
program\_2.s       & 0.01778 & 0.04457 & 0.00446 & 0.02674 \\
program\_3.s       & 0.09646 & 0.06270 & 0.24114 & 0.00916 \\
\textbf{TOTAL Time (@ 1.75kHz)} & \textbf{2.12143} & \textbf{1.53523} & \textbf{2.41354} & \textbf{2.05525} \\
\hline
\end{tabularx}
\end{table}

\end{document}

